\documentclass{article}
\usepackage{amsmath,amssymb,amsthm}
\newtheorem{proposition}{Proposition}
\title{An exact initial seam-variance bound for the original Wilson gluing}\author{}\date{}
\begin{document}\maketitle
We use precisely the same-variable gluing in manuscript Section 92, equations SV1--SV12. The group is $SU(N)$, $N\geq2$, links have normalized Haar measure, and $s(U)=1-N^{-1}\operatorname{Re}\operatorname{tr}U$. The same Wilson coefficient $\beta$ is used in both systems. This note proves a new finite-cutoff variance statement; it does not prove the signed integral required by S94.

Write $n=\ell_*/a\geq3$, $q\geq3$, $M=q^4$, and $\mu_0=\mu_{a,T_*}^{\otimes M}$. Every link retains its original label $(b,x,\mu)$, where $b\in(\mathbb Z/q\mathbb Z)^4$, $x\in\{0,\ldots,n-1\}^4$. Set $\Delta S=S_{\rm big}-S_{\rm per}$ as in SV5. No gauge is fixed.

\section{Orthogonality under the original product gauge group}
The probability $\mu_0$ is invariant under independent site gauge transformations in each small periodic torus. Call a new big-torus plaquette \emph{singly rewired} if exactly one of its two directional coordinates is $n-1$. The exact number is
\[
 N_1=12M(n-1)n^2.
 \tag{1}
\]
Indeed each positive orientation has $2(n-1)n^2$ such starts, and there are six orientations. This count includes each plaquette once.

Suppose the crossed direction is $\mu$ and the tangential direction is $\nu$. Put $y=x+e_\mu$ reduced within the small block, so $y_\mu=0$. After the factors which cancel within the old link products are canceled, the new plaquette trace has four unmatched original gauge vertices:
\[
 (b,y),\quad(b,y+e_\nu),\quad
 (b+e_\mu,y),\quad(b+e_\mu,y+e_\nu).
 \tag{2}
\]
Each occurs once in the matrix word, in the fundamental or its dual representation. Integrating any one of these gauge variables annihilates the trace. Thus for
$X_P=N^{-1}\operatorname{Re}\operatorname{tr}U_P^{\rm big}$,
\[
 E_0X_P=0.
\]
Holding all other gauge transformations and links fixed, that same gauge variable makes the plaquette matrix Haar-distributed, up to fixed left and right unitary factors. Elementary fundamental-representation Haar orthogonality gives
\[
 E_0X_P^2=\kappa_N,\qquad
 \kappa_N=\begin{cases}1/4,&N=2,\\1/(2N^2),&N\geq3.\end{cases}
 \tag{3}
\]
For completeness, $\int|\operatorname{tr}U|^2dU=1$. Also $\int(\operatorname{tr}U)^2dU=0$ for $N\geq3$, by the nontrivial center character (including $N=3$), whereas for $SU(2)$ the trace is real. These prove (3).

For distinct singly rewired plaquettes the four-vertex sets (2) are different when $q,n\geq3$. The adjacent block pair determines the crossed direction and initial block, and the unordered tangential nearest-neighbor pair determines the tangential direction and start. The restrictions $q,n\geq3$ exclude the two-cycle ambiguity. A vertex belonging to one set and not the other appears only once in their product, so gauge averaging gives
\[
 E_0X_PX_Q=0\quad(P\ne Q,\ P,Q\hbox{ singly rewired}).
 \tag{4}
\]
A doubly rewired plaquette has six unmatched gauge vertices in three small blocks: two in $b$, two in $b+e_\mu$, and two in $b+e_\nu$. These are distinct for $q,n\geq3$. Its vertex set differs from (2); if the four-vertex set were contained in it, one of its extra vertices could instead be integrated. Thus every singly rewired $X_P$ is orthogonal to every doubly rewired trace as well. It is also orthogonal to every old plaquette action or any other old gauge-invariant integrable function.

\section{An actual Wilson variance lower bound}
\begin{proposition}
For every $\beta\geq0$ and every $q,n\geq3$,
\[
 \operatorname{Var}_{\mu_0}(\Delta S)\geq
 \kappa_N\,12M(n-1)n^2.
 \tag{5}
\]
\end{proposition}
\begin{proof}
Let $Y=\sum_{P\ {
m singly\ rewired}}X_P$. In $\Delta S$ every such trace has coefficient $-1$. All remaining new traces and all old action terms are orthogonal to $Y$, by the gauge averages above. Constants are orthogonal as well since $E_0Y=0$. Hence $\Delta S=-Y+R$ with $\operatorname{Cov}_0(Y,R)=0$. Equations (3)--(4) give
$\operatorname{Var}_0(\Delta S)=\kappa_NN_1+\operatorname{Var}_0(R)$.
\end{proof}
This is a same-Wilson statement for the literal gluing, with no weak-field approximation, selector modification, or new normalization. It does not assert that the variance remains this large away from $t=0$.

\section{The exact cancellation scale required by S94}
On the same compact Haar space let
$d\mu_t=Z_t^{-1}e^{-\beta(S_{\rm per}+t\Delta S)}dU$.
Differentiation under the finite integral gives
\[
 \frac{d}{dt}E_t\Delta S=-\beta\operatorname{Var}_t(\Delta S),
\]
and Taylor's integral formula for $\log Z_t$ yields
\[
 G:=M\log Z_{a,T_*}-\log Z_{a,qT_*}
 =\beta E_0\Delta S-
 \beta^2\int_0^1(1-t)\operatorname{Var}_t(\Delta S)\,dt.
 \tag{6}
\]
Here $E_0\Delta S=N_{\rm seam}(1-\bar s)$ with
$N_{\rm seam}=M(12n^3-6n^2)$ and the original small-torus plaquette expectation $\bar s$. Consequently the already available $G\geq0$ implies
\[
 \int_0^1(1-t)\operatorname{Var}_t(\Delta S)\,dt
 \leq \frac{N_{\rm seam}(1-\bar s)}{\beta}.
 \tag{7}
\]
Conversely the desired uniform upper bound $G\leq C M\ell_*^4$ is equivalent, for $\beta>0$, to
\[
 \int_0^1(1-t)\operatorname{Var}_t(\Delta S)\,dt
 \geq \frac{N_{\rm seam}(1-\bar s)}{\beta}
 -\frac{CM\ell_*^4}{\beta^2}.
 \tag{8}
\]
Equation (5) identifies an initial variance of order $Mn^3$, whereas (7) bounds its weighted integrated mass by order $Mn^3/\beta$. Thus at weak coupling its initial size cannot persist on a fixed interpolation interval. This is an exact restriction on any proposed uniform-in-$t$ approximation. It is not a proof of (8): a value at one endpoint does not control the required integrated cancellation. In particular, replacing the seam energy by a local small-plaquette estimate under $\mu_0$ fails both at the mean and at the variance level.

All independent absolute scalar factors of the original RG construction remain outside this normalized expectation calculation. The relation to its vacuum $V$ uses the original stability comparison SV3 only. No absolute scalar increment or inherited carrier has been canceled by (5)--(8). Full cross-torus S94 remains open.
\end{document}
